% !TeX spellcheck = ru_RU
% !TEX root = vkr.tex
\section{Постановка задачи}
\label{sec:task}

Цель работы --- развить рекомендательную систему PySATL-expert для поддержки
выбора вероятностного распределения по одномерной числовой выборке, обеспечив
настраиваемый статистический анализ и обоснованное представление рекомендаций.

Для достижения цели поставлены следующие задачи:
\begin{enumerate}
    \item Развить архитектуру PySATL-expert, обеспечив настройку статистического
    анализа и расширение набора поддерживаемых распределений, критериев
    и моделей рекомендаций.
    \item Разработать способы представления результатов статистических
    проверок для формирования рекомендаций с учётом оценивания неизвестных
    параметров.
    \item Исследовать методы формирования рекомендаций и экспериментально
    сравнить их по качеству и вычислительным затратам.
\end{enumerate}

Вход системы включает выборку и конфигурацию анализа. Истинное семейство
и параметры генератора, известные в модельном эксперименте, анализатору
не передаются. Результат включает рейтинг кандидатов, оценки их параметров,
результаты выполненных проверок и сведения об ограничениях вычислений.
Оценка модели не интерпретируется как вероятность истинности
статистической гипотезы.

Промежуточный результат до декабря --- работающий расширенный конвейер
на выбранном наборе распределений и критериев и первый сравнительный
эксперимент по представлениям результатов проверок при одинаковых модели
и разбиении данных. Исследование методов рекомендаций и основное сравнение
продолжаются во втором семестре.

Обзор используется для обоснования требований и выбора исследуемых
подходов. Преимущество нового представления или способа обучения
не предполагается заранее и должно устанавливаться экспериментально.
